% ============================================================
%   §0 Abstract
% ============================================================
I analyse a high-resolution spatial transcriptomics cohort of $\sim 8$ million individually segmented cells from 56 colorectal-cancer patients (112 samples), annotated with eight tumour-microenvironment cell types: B cells, Endothelial cells, Fibroblasts, Myeloid cells, Normal epithelium, SMC, T cells, and Tumor cells. The data are intrinsically network-shaped --- every cell has spatial neighbours that condition its biology --- and I treat the spatial $k$-nearest-neighbour graph as the integrating object across a five-Part analysis pipeline.

\partref{1} characterises the topology of the spatial graph across 110 non-anomaly samples. Local statistics are remarkably consistent (median mean degree $8.24$), but the global fraction of cells in the largest connected component spans a continuous gradient from $0.11$ to $0.999$, from compact to fragmented tissue. Topology correlates with cell-type composition: SMC drives sparsity ($\rho = -0.58$ vs.\ mean degree), and per-cell betweenness identifies a tissue skeleton with no cell-type signature.

\partref{2} aligns per-sample Leiden communities into 20 stable niches across the cohort, covering 77\% of cohort cells. Niches partition into four recurrence bands: 4 universal, 11 common, 4 selective, and one rare (a pure-SMC niche restricted to 11 patients). The niches map asymmetrically onto the eight cell types: Tumor cells, B cells, and Fibroblasts each dominate four niches; T cells dominate none. The only niche where T cells appear as a substantial secondary minority is \nicheid{17} (48\% B cells, 29\% T cells), a composition consistent with tertiary lymphoid structures.

\partref{3} runs a five-rung ablation comparing per-cell baselines (logistic regression, multi-modal MLP) to graph-aware architectures (GraphSAGE, GAT) on the cell-type classification task. A 2-layer GraphSAGE on highly variable gene features alone achieves the best validation macro-$F_1$ ($0.6350$), a $+6.5$pp gain over the strongest per-cell baseline ($0.5704$), and a final test macro-$F_1$ of $0.6089$ on held-out patients. Adding the H\&E embedding or graph attention on top of GraphSAGE yields no further improvement: the spatial graph subsumes the architectural signal that the H\&E modality contributes in isolation.

\partref{4} fits per-class Gaussian Graphical Models on 378 broadly variable common genes. At the near single-cell resolution of Visium HD, mRNA diffusion between adjacent cells means that per-cell expression vectors carry a non-trivial component of neighbourhood signal, and the resulting networks recover \emph{microenvironmental gene programmes} rather than purely cell-intrinsic regulation. Five of the eight classes retain canonical cell-type markers among their top hubs alongside neighbourhood-derived genes; three (B cells, Endothelial cells, T cells) are dominated by neighbourhood signal. Across the eight networks, hub Jaccard reaches $0.36$ while edge Jaccard stays below $0.09$: cell types share hubs but rewire connections.

\partref{5} repeats the GGM fit per-niche rather than per-class. Pairwise edge Jaccard between niches correlates moderately with composition cosine (Spearman $\rho = 0.615$), but the cluster-level family structure of \partref{2} is not recovered (Adjusted Rand Index $0.25$ for edges, $0.05$ for hubs, vs.\ $0.80$ for composition itself). Patient cohort overlap explains part of the structure ($\rho = 0.42$) but composition is the stronger driver. \nicheid{17}, the B-cell + T-cell niche, is the least edge-integrated of the B-cell-dominant niches --- a weak signal on an essentially flat dendrogram rather than a clean outlier --- in a direction consistent with the distinctive microenvironmental gene programme of T cells identified in \partref{4}.

The most consistent finding across all five Parts is the structural distinctness of T cells, flagged by four independent analyses (negative cell-type assortativity in \partref{1}, niche absence in \partref{2}, classification difficulty in \partref{3}, gene-network disjointness in \partref{4}). The unifying explanation is biological: T cells in tumour tissue are diffuse infiltrating populations with no consistent spatial neighbourhood, except in the TLS-like \nicheid{17} where they coexist with B cells in organised aggregates. The most generalisable methodological finding is the dominance of the spatial graph over per-cell modalities for the supervised task, and the convergence of cell-type and niche-level gene networks on a common "shared hubs, rewired connections" architecture.